%! Fundamental Theorem of Algebra and Complex Zeros — Unit 4, Lesson 4.5 — Homework
\documentclass[10pt]{article}
\usepackage{algebra2-article}
\usepackage{algebra2-boxes}
\newcommand{\TallMath}[1]{$\displaystyle #1\rule[-1.4em]{0pt}{3.2em}$}

\begin{document}

\pageheader{Unit 4, Lesson 4.5}{Homework}
\namedateperiod

\vspace{0.10in}

\begin{notesbox}{Practice}
The degree tells you how many zeros to find. Conjugate pairs tell you that imaginary zeros never come
alone. Count every answer against the degree before you move on.

\begin{enumerate}[leftmargin=1.9em, itemsep=8pt]
  \item \textbf{Just count.} Each polynomial is already factored. Give the degree and the zeros
    counted with multiplicity.
    \begin{enumerate}[label=(\alph*), leftmargin=1.8em, itemsep=5pt]
      \item $f(x) = (x-7)(x+2)^3$: \; degree \blank{1.0cm}, zeros \blank{4.4cm}
      \item $g(x) = (x^2+1)(x-5)^2$: \; degree \blank{1.0cm}, zeros \blank{4.4cm}
      \item $h(x) = (x^2+4)^2(x-1)$: \; degree \blank{1.0cm}, zeros \blank{4.4cm}

        \vspace{2pt}
        In (c), how many of those zeros are imaginary? \blank{1.2cm} \; Is that number even?
        \blank{1.0cm}
    \end{enumerate}

  \item \textbf{Number and type, without solving.} Each polynomial has \emph{real} coefficients.
    \begin{enumerate}[label=(\alph*), leftmargin=1.8em, itemsep=5pt]
      \item Degree $3$; its graph passes straight through the $x$-axis exactly once. \;
        Real: \blank{1.0cm} \; Imaginary: \blank{1.0cm}
      \item Degree $4$; its graph \emph{touches} the $x$-axis at $x=3$ and meets it nowhere else. \;
        Real: \blank{1.0cm} \; Imaginary: \blank{1.0cm}
      \item Degree $5$; its graph passes straight through at $x=0$, $x=2$, and $x=-4$, and nowhere
        else. \; Real: \blank{1.0cm} \; Imaginary: \blank{1.0cm}
      \item Degree $6$; its graph never touches the $x$-axis. \; Real: \blank{1.0cm} \;
        Imaginary: \blank{1.0cm}
    \end{enumerate}

  \item \textbf{Solve over $\mathbb{C}$.} $x^3 + 2x^2 + 4x + 8 = 0$ \; (this is the quotient
    Lesson 4.3's homework produced when you divided $x^4-16$ by $x-2$ --- now finish it).
    \begin{enumerate}[label=(\alph*), leftmargin=1.8em, itemsep=5pt]
      \item Factored form (grouping works): \blank{4.6cm}
      \item All solutions: \blank{4.0cm}
      \item Number and type: \blank{1.0cm} real, \blank{1.0cm} imaginary
    \end{enumerate}

  \item \textbf{Solve over $\mathbb{C}$ --- this one needs the full workflow.}
    $x^3 - 7x^2 + 17x - 15 = 0$
    \begin{enumerate}[label=(\alph*), leftmargin=1.8em, itemsep=5pt]
      \item Candidates: \blank{5.0cm} \; One that works: $r = $ \blank{1.0cm}
      \item Depressed polynomial: \blank{2.8cm} \; Its discriminant: \blank{1.6cm}
      \item All solutions: \blank{4.6cm}
      \item \textbf{Verify} one of the imaginary solutions is paired correctly, and explain how the
        degree confirms you are finished. \blank{6.0cm}
    \end{enumerate}

  \item \textbf{Closing a loop.} Solve $x^4 - 16 = 0$ completely.
    \begin{enumerate}[label=(\alph*), leftmargin=1.8em, itemsep=5pt]
      \item Factor as a difference of squares --- \emph{twice} if you can: \blank{4.6cm}
      \item All four solutions: \blank{4.0cm}
      \item In Lesson 4.2 you would have said $x^2+4$ is prime and stopped. Is it still prime over
        the \emph{complex} numbers? \blank{1.2cm} \; Factor it: \blank{2.6cm}
    \end{enumerate}

  \item \textbf{Build them.} Write each polynomial in factored form \emph{and} standard form.
    \begin{enumerate}[label=(\alph*), leftmargin=1.8em, itemsep=5pt]
      \item Zeros $3$ (multiplicity $2$) and $-1$. \; Factored: \blank{3.4cm} \;
        Standard: \blank{4.2cm}
      \item Real coefficients; zeros include $-4$ and $2i$. \; Full zero list: \blank{3.0cm}

        \vspace{2pt}
        Factored: \blank{3.4cm} \; Standard: \blank{4.6cm}
    \end{enumerate}

  \item \textbf{Deduce the missing one.} A polynomial of degree $5$ with real coefficients has zeros
    $3$, $-1$ (multiplicity $2$), and $2-i$.
    \begin{enumerate}[label=(\alph*), leftmargin=1.8em, itemsep=5pt]
      \item How many zeros have you been given, counted with multiplicity? \blank{1.2cm}
      \item What is the fifth zero, and how do you know? \blank{5.4cm}
      \item How many $x$-intercepts does the graph of this function have? \blank{1.2cm}
    \end{enumerate}

  \item \textbf{Explain.} Answer each in one or two sentences.
    \begin{enumerate}[label=(\alph*), leftmargin=1.8em, itemsep=5pt]
      \item Why must every polynomial of \emph{odd} degree with real coefficients have at least one
        real zero? Give both reasons --- the conjugate-pair one and the end-behavior one.
        \writelines{2}
      \item A classmate solves a cubic, finds one real solution, and says ``the other two are
        missing.'' Rewrite their sentence so it is correct, and say what a graph can and cannot show.
        \writelines{2}
    \end{enumerate}
\end{enumerate}
\end{notesbox}

\begin{extensionbox}
\textbf{Design and justify.}
\begin{enumerate}[label=(\alph*), leftmargin=1.8em, itemsep=5pt]
  \item Show that for any real numbers $a$ and $b$, the product
        $\big(x-(a+bi)\big)\big(x-(a-bi)\big)$ has real coefficients, and state what it equals.
        (Group the factors as $(x-a) \mp bi$ first.)
  \item What is the smallest possible degree of a polynomial with real coefficients whose zeros
        include $0$, $\sqrt{5}$, and $3-2i$? Explain which zeros do and do not force a partner ---
        and be careful, because only one of these rules applies to $\sqrt{5}$.
  \item Write a quartic with real coefficients whose graph has \emph{exactly two} $x$-intercepts,
        then a quartic whose graph has \emph{none}. Give each in factored form and justify the number
        of intercepts using the counting rules, not a graph.
  \item Explain why no polynomial with real coefficients and degree $4$ can have exactly \emph{three}
        $x$-intercepts, each a simple crossing, and no other zeros.
\end{enumerate}
\writelines{3}
\end{extensionbox}

\begin{spiralbox}
\textbf{Coming next --- Lesson 4.6: Graphing Polynomial Functions \& Modeling.} You now own every
piece of a polynomial's graph separately: end behavior from the degree and leading coefficient
(4.0), zeros from factoring and the Rational Root Theorem (4.2--4.4), how the graph behaves at each
zero from its multiplicity (4.0), and --- as of today --- exactly how many zeros exist and how many of
them the graph is allowed to show. Tomorrow those pieces become one procedure: read a polynomial,
predict its graph, and use it to answer questions about a real situation.
\end{spiralbox}

\end{document}
